% ================================
% \section*{\centering \bf Part V: Future and Challenges}
% \addcontentsline{toc}{section}{Part V: Future and Challenges}
% ================================

\begin{figure}
    \centering
    \includegraphics[width=\linewidth]{figures/part5_opener_v2.pdf}
    
    \caption{Part V in one view. (1) What it requires: a household robot must keep one consistent account of its kitchen while it turns away, people move things, and a stack of cans may topple. (2) Where the gaps are: the loop from perception to action still breaks, benchmark success diverges from deployment, and consistency erodes over long horizons. (3) What comes next: open problems that set the directions for future work.}
    \label{fig:p5_main}
\end{figure}

Parts~II to~IV described how spatial representations are constructed, how they support perception, reasoning, and generation, and how they are deployed in applications ranging from household robots and surgical systems to city-scale digital twins.
However, current systems often perceive or act well at a single moment, while most of the open problems arise when a system must maintain a consistent picture of the world as it moves, revisits places, and acts over extended periods.

A household robot that puts away groceries must remember the location of a cupboard after turning away from it, register a chair that has been pushed into its path, track which bags it has already emptied, and judge whether a stack of cans will remain stable if another can is added.
Methods described earlier in this handbook handle each of these steps reasonably well in isolation.
The difficulty lies in maintaining a single account of the kitchen that all steps read from and write to, and in keeping that account correct while the robot and the people around it change the scene.
In the formalism of Chapter~\ref{chapter:definition}, this account is the world state $\mathcal{R}^{\mathcal{S}}_t$, and keeping it correct is the role of the update $\Psi_{\mathcal{S}}$.
It is also the problem that Section~\ref{sec:historical-evolution} identifies as open for the current era, long-horizon consistency under motion and partial observability.

Chapter~\ref{chapter:p5_discussions} asks what makes such a world state usable: what can be inferred from limited observations, what must be represented, how beliefs are kept current, when predictions can be trusted, and how spatial knowledge is acquired through interaction in shared space.
We present the evidence on each question and state which answers remain open.
Chapter~\ref{chapter:p5_challenges} turns to gaps in current methods and benchmarks, the difficulties of deployment, and directions for future work. 
Chapter~\ref{chapter:p5_conclusion} concludes by revisiting the definition of spatial intelligence introduced in Chapter~\ref{chapter:definition}.

\newpage

\chapter{Foundations: What Spatial Intelligence Requires}
\label{chapter:p5_discussions}



% \lettrine[lines=3]{\initfamily\textcolor{darkgreen}{O}}{pen}
% Vision-based ranger involves leveraging geometric principles to calculate distance.
% A vision-based ranger could use one or multiple cameras.

\FramedInitial{A}{dvances} in spatial perception, reasoning, and generation have enabled AI systems to recover geometry from images, answer questions about layout, and simulate complex environments, as Parts~II to~IV documented.
Success on these tasks does not establish that a system maintains a reliable model of the physical world.
A model can reconstruct a room to millimeter accuracy, answer most questions about a photograph, or continue a video convincingly, and still fail when observations are incomplete, when the environment changes while it is not looking, or when its actions depart from the conditions under which it was trained.
The failures reported throughout this handbook cluster at these points, which suggests a shared deficiency rather than a collection of unrelated task-level shortcomings.

This chapter asks what makes an internal model of physical space usable for continued reasoning and action.
By usable we mean that the model can be queried for what a task requires, extended when new evidence arrives, and relied upon when its predictions are acted on; we do not presuppose that such a model must be an explicit 3D reconstruction, nor that it must be a video generator.
The question decomposes into five, each posed at a different level.
What can be inferred from limited observations, and what remains underdetermined (Section~\ref{sec:p5_observability})?
Which information must a representation retain, and for which operations (Section~\ref{sec:p5_sufficiency})?
How can a world state be maintained and revised over time (Section~\ref{sec:p5_persistence})?
When can its predictions be trusted (Section~\ref{sec:p5_validity})?
How is spatial knowledge acquired and used through interaction, in environments shared with other agents (Section~\ref{sec:p5_agency})?
The five are interdependent: action produces new evidence, which re-enters at the first level and constrains all of the others.


\section{Spatial Observability and Identifiability}
\label{sec:p5_observability}

\FramedInitial{S}{patial} understanding begins with an incomplete view of the world.
A system never has direct access to the world state $\mathcal{R}^{\mathcal{S}}_t$ of Chapter~\ref{chapter:definition}; it has images, video, depth, or sound, from which that state must be inferred.
Images and videos carry rich evidence about appearance and structure, but they do not uniquely determine every geometric, physical, or functional property of a scene.
A single image fixes the scale of a scene only through the prior knowledge brought to it (Section~\ref{sec:height_distance}), an occluding surface hides whatever lies behind it, and the mass of a box or the state of a lock leaves no trace in a photograph.
A usable model must therefore keep track of which of its contents are observed and which are inferred.

Whether a spatial inference fails because the model cannot extract the evidence available, or because the evidence itself is insufficient, is a distinction current evaluation rarely draws.
The former is a limitation of learning or computation, and better models or more data can be expected to reduce it.
The latter is a problem of identifiability, which more capable inference alone does not remove.
Asked whether a sofa seen in one photograph will fit through a door seen in another, a multimodal model typically produces a confident answer with a coherent explanation, yet without a common scale between the two images the question has no determinate answer, and a model with perfect geometric reasoning should say so.
Evaluation that counts every wrong answer as a model failure conceals this class of cases, while evaluation that accepts every ambiguity as irreducible would excuse failures that are not.

Additional modalities constrain different aspects of the underlying state, and their contributions are neither interchangeable nor universally sufficient.
A second viewpoint supplies parallax, which fixes relative depth and reveals part of what the first view occluded.
A depth sensor supplies distances directly, with an accuracy and a scale behavior that depend on the sensing principle (Section~\ref{sec:depth_estimation}).
Motion adds constraints on dynamic geometry; sound carries direction, distance, and room structure (Section~\ref{sec:av_localization}); contact and force reveal mass, friction, and compliance, which no camera observes (Section~\ref{subsec:p3_s7_1_2}).
Assessing a modality therefore requires specifying which ambiguities it resolves, under which assumptions, and which properties remain unidentifiable once it has been added.
Part~II treated the modalities as alternative representations of space (Chapter~\ref{chapter:single_modality}); here they appear as complementary constraints on one underdetermined state.

Observability is not fixed by the initial input.
Moving to another viewpoint, changing the illumination, or pushing an object can reveal previously ambiguous properties, and some properties, such as a drawer's operability, may require physical interaction.
Theory of Space and VABench bear directly on this point: models perform better with supplied observations than when they must gather evidence themselves, and allowing a model to move its camera and look again raises its manipulation success substantially \cite{zhang2026theory,zhang2026vabench}.
Spatial intelligence therefore depends partly on recognizing when an additional observation would change what is knowable.
How an agent should choose such observations is a question of decision making, taken up in Section~\ref{sec:p5_agency}, which properties are identifiable depends on the agent's available actions.

Distinguishing uncertainty from fundamental underdetermination, and stating what evidence would resolve it, are abilities that current systems are not asked to demonstrate.
Static question accuracy cannot distinguish correct reasoning under insufficient evidence from incorrect reasoning despite sufficient evidence, and controls as simple as removing the image or shuffling the input remain rare (Section~\ref{sec:vl_grounding}).
A benchmark designed around identifiability would include questions that are unanswerable from the given input and reward a model for saying so, pair each ambiguous case with the modality or action that resolves it, and score whether the model asks for it.
Given these limits, which obtainable properties should a spatial representation preserve?


\section{Representation and State Sufficiency}
\label{sec:p5_sufficiency}
\label{sec:p5_cognitive}

\FramedInitial{T}{he} same physical environment can be represented through images, point clouds, occupancy grids, scene graphs, neural fields, or the latent state of a learned model, and Part~II surveyed each of these forms (Chapter~\ref{chapter:single_modality}).
They differ less in how much information they hold than in which aspects of the environment they preserve, which they abstract away, and how they organize what remains.
Scene graphs typically emphasize entities and their relations, Gaussian representations primarily encode appearance and geometry, and a latent state retains what its training objective rewarded; each can be extended with further structure, but none is informative in isolation from the questions and operations it is meant to support.

That spatial knowledge and spatial reasoning are separable is shown by large models that acquire the former without the latter.
A transformer trained only on Othello move sequences forms an internal representation of the board (Section~\ref{sec:rep_language}), and language models encode the geographic positions of cities in their activations without ever being shown a map \cite{gurnee2023language}.
The same models remain well below human performance on perspective taking, mental rotation, and building a model of a place from a sequence of views \cite{cai2025holistic}, and chain-of-thought prompting, which helps on arithmetic and code, tends to degrade their accuracy on visual spatial benchmarks \cite{kancheti2026cotdegrades}.
Performance recovers when the model is allowed to draw a coarse cognitive map, render an imagined view, or call a geometric tool before answering (Sections~\ref{subsec:p3_s7_3_5} and~\ref{subsec:p3_s7_3_8}).
Encoding spatial information and computing with it are therefore distinct requirements, and reasoning performance depends partly on whether the representation in use supports the operations the task demands.
This does not show that language is unsuited to spatial reasoning; it suggests that a representation in which distances and viewpoints are described rather than held tends to accumulate error each time they are manipulated.

Spatial tasks differ not only in the information they require but in the computations they demand of it.
Judging a relative position needs a reference frame and a relation; metric reasoning needs distances, angles, and a scale; changing viewpoint needs a rigid or projective transformation applied to 3D structure; navigation needs topological connectivity and traversability; manipulation needs local geometry together with contact constraints and material properties.
A representation suited to semantic relations can be insufficient for metric transformation, and a detailed geometric reconstruction can still omit the properties on which physical interaction depends, a gap Section~\ref{sec:p4_s11_1} summarized as geometry is not behavior.
The effectiveness of a representation is therefore a matter of its compatibility with the downstream operations.
The recurring finding of Part~III that hybrid systems are limited by their weakest interface (Section~\ref{sec:p3_s7_4}) is an instance of this: the interface is where one representation's operations are asked to run on another representation's content.

These observations suggest a task-relative criterion of sufficiency.
Spatial-state sufficiency depends on the queries, predictions, and actions it must support; a representation is sufficient for a family of tasks if every decision the family requires can be made from it as well as from the world itself.
Information irrelevant to one task may be essential to another, so sufficiency cannot be judged from the representation alone.
This creates a tension between compact states optimized for an immediate objective and general-purpose states intended to remain useful across tasks that are not yet known.
A latent world model can succeed at the tasks it was trained for while modeling the environment's dynamics unfaithfully \cite{gonzalezlaiz2026limits}, which satisfies sufficiency for those tasks and gives no assurance for others; a dense geometric reconstruction, by contrast, serves many queries but remains silent about physical properties that are not visible.

Whether a finite shared state can preserve the distinctions required by an open-ended range of tasks, or whether a spatial system should instead maintain adaptable representations whose detail and structure depend on downstream demands, remains open.
Shared spatial latents read by many task heads, task-specific representations maintained in parallel, and interfaces through which modules request information at the precision they need (Section~\ref{sec:p3_s7_4}) represent three answers, and the available evidence does not yet favor one, particularly when future tasks are unknown.

Evaluating a spatial representation therefore requires examining the information it preserves and the operations it supports, rather than its expressiveness in the abstract.
Reconstruction quality and recognition accuracy are the usual proxies, and both are insufficient: a photometrically accurate replica can fail a metric tolerance check and encode nothing about which drawers open (Section~\ref{sec:p4_s11_1}), and a correct answer can rest on language priors rather than on a correct belief about the scene.
Sufficiency has to be tested across tasks, by asking whether a representation fixed for one purpose supports queries, transformations, and action predictions it was not built for.
Whether a state that is sufficient at one moment remains so as the world changes and objects leave the field of view is a separate question, taken up next.


\section{Persistent World States and Belief Updating}
\label{sec:p5_persistence}
\label{sec:p5_worldstate}

\FramedInitial{U}{nderstanding} a scene at one moment differs from maintaining an understanding of a changing world.
An assistant asked where the keys were left must have observed them, retained their location after the user walked away, distinguished that remembered location from the current view, and allowed for the possibility that someone moved them in the meantime.
Egocentric benchmarks built around such queries show that an explicit memory, updated as video arrives and tracking objects relative to the moving user, substantially improves leading multimodal models \cite{wang2026keepitinmind}, whereas most commercial glasses answer from a single frame and keep no map (Section~\ref{sec:p4_s9_3}).
The household robot in this part's introduction faces the same requirement at every step.
In the formalism of Chapter~\ref{chapter:definition} this is the role of the update $\Psi_{\mathcal{S}}$, the component that Parts~II to~IV examined least.

Access to a long observation history does not ensure a consistent world state.
Video world models steered through long explorations fail to render revisited places consistently even as image quality improves \cite{wang2026worldprobe}; multimodal models given several views of a room rarely detect that one view was captured after an object moved \cite{cao2026samespace}; and over routes spanning several streets, models answer by recognizing landmarks and matching text, with little evidence of path integration \cite{mahendran2026kilometervision}.
A model may retain every earlier frame and still fail to recognize a revisited place, preserve an object's identity, or decide which remembered properties have become outdated.
Research on spatial memory has concentrated on recall, through longer contexts, better retrieval, and richer memory structures (Section~\ref{subsec:p3_s7_3_6}), while these failures concern revision.
Persistence therefore involves maintaining and revising beliefs about the world, not merely retaining observations.

A persistent world state must distinguish absence of observation from evidence of absence, maintain uncertainty during occlusion, and reconcile new observations with prior beliefs.
When a new observation conflicts with a stored one, the system must decide whether the environment changed, the earlier estimate was wrong, or the two observations refer to different entities, each requiring different state updates.
Classical systems partly address this problem: SLAM detects loop closures and corrects the map upon revisiting a known place (Section~\ref{sec:mapping_slam}), object trackers maintain identities through occlusion, and dynamic scene graphs timestamp their edges.
Learned systems often struggle to make such revisions reliably over extended interactions.
Models retain outdated beliefs after the world changes \cite{zhang2026theory}, append-only memories accumulate contradictory entries, and agents drift over long episodes as their context fills (Section~\ref{subsec:p3_s7_3_7}).
Increasing capacity alone may not suffice, since a memory that only grows keeps each outdated belief beside its replacement.
A write policy that decides what enters the state, with what confidence, and when an entry is revised or removed, is a prerequisite that current systems leave largely implicit.

Where such a state should be kept is a question of architecture rather than of representation format.
Explicit structures such as maps, scene graphs, and splats (Section~\ref{sec:rep_3d}) persist by construction and make selected state variables accessible and correctable, but they must anticipate every type of object, property, and change they will be asked to encode.
Learned memories and recurrent latent states can encode properties no designer specified, and whether they track the world is difficult to verify \cite{gonzalezlaiz2026limits}.
Generated frames are the most flexible option and lose whatever leaves the context window.
Hybrid designs that store static geometry and dynamic objects separately and retrieve them on revisit \cite{yu2026grome}, or that keep slowly changing layout and identities in an explicit layer and leave appearance and short-term dynamics to a learned one, combine these advantages at a cost of their own: geometry, object identity, uncertainty, and temporal updates must stay synchronized across layers, and coordinate frames and scales can silently diverge at the interface (Section~\ref{sec:p3_s7_4}).
The separation can be taken further.
State World Models keep a typed state ledger that evolves under executable rules written by a coding agent and use a frozen video model only to render each resulting event; on counterfactual twin pairs that share an initial frame and action but differ in hidden state, this harness produces the rule-required outcome in 92\% of pairs, where the best of eleven video models prompted with the same rules reaches 32\% \cite{chan2026stateworld}.
The result suggests that what fails in generator-only designs is the maintenance of state the current frame does not show, rather than rendering.
The appropriate architecture follows from the requirements of state maintenance, how long each part must persist and whether it must be inspected, rather than from a general preference for explicit or learned form.

World states should remain persistent without becoming rigid and revisable without losing established structure, and evaluation currently measures neither property.
Most benchmarks refer to a single moment, and the protocols that would expose persistence failures are not yet standard (Section~\ref{sec:p3_s7_5}): streaming protocols that withhold future frames \cite{li2026ovosbench}, revisit protocols that move the camera away and back \cite{wang2026worldprobe}, and change protocols that alter the scene behind the agent.
A more informative evaluation reads out the state itself, as Theory of Space does at each exploration step \cite{zhang2026theory}, and scores separately what the system has observed, what it remembers, and what it predicts, since conflating the three is how a stale belief passes as a current one.
Deployed systems could likewise report how far their world state lags behind the world it describes, a quantity that digital twins call the liveness gap (Section~\ref{sec:p4_s11_1}) and that is rarely measured elsewhere.
A state that remains consistent over hours can nonetheless describe appearance accurately and physical behavior poorly.


\section{Predictive Validity and Physical Grounding}
\label{sec:p5_validity}
\label{sec:p5_plausible}

\FramedInitial{A}{} video generator asked to continue a clip of a ball rolling off a table usually produces a smooth, well-lit continuation in which the ball falls, often along the wrong trajectory or with an implausible bounce.
Benchmarks of physical plausibility find that the physical accuracy of video generators is poor and does not improve with visual quality (Section~\ref{sec:p3_s7_5}), and Chapter~\ref{sec:generation} reported the same pattern across its subfields: shapes that change between views, feet that slide over the ground, and motions that violate conservation (Section~\ref{sec:p3_s8_8}).
A plausible continuation reproduces the appearance of a physical event without predicting its trajectory, contacts, or consequences.
Yet the same class of models learns substantial scene structure, since image generators adapted for perception segment, estimate depth, and ground referring expressions at the level of specialist models \cite{gabeur2026image}.
Structure and dynamics are apparently learned to different degrees from the same data.

World-model correctness has several distinct meanings, and the criteria are often run together.
Visual plausibility concerns realistic appearance; geometric consistency concerns coherent spatial structure across views and time; dynamical fidelity concerns how the world evolves under actions; decision reliability concerns whether actions planned in the model succeed in reality.
The criteria are related but not interchangeable, and they are ordered by how much they depend on action.
Common video metrics measure the first and do not register physically impossible motion (Section~\ref{subsec:p3_s8_3_6}); view-synthesis metrics measure the second and cannot tell a solid object from a cut-out that matches the training views (Section~\ref{subsec:p3_s8_2_1}); the third has no standard protocol (Section~\ref{subsec:p3_s8_6_discussion}); the fourth is tested, when it is tested at all, by deploying the resulting policy.
A model can score well on the first two and fail the last two, and vision-language models used as judges share the spatial blind spots of the models they evaluate (Section~\ref{sec:p3_s8_8}).

Learning how scenes typically evolve from observational data does not guarantee correct predictions under novel interventions.
In the action-conditioned definition of Chapter~\ref{chapter:definition}, $f^{\mathcal{S}}_{\text{pred}}$ maps a state and an action to the next state, but a generator trained on video sees only the actions that happened to be recorded, mostly everyday activity and camera motion, and sees them correlated with the states and circumstances in which they occurred.
A model that has watched thousands of doors being pushed open has learned what usually follows a push, not what determines the outcome, and so has little basis for predicting the locked door, the obstructed door, or the push applied at an unusual angle.
The regularities it learns are real and useful for geometry, appearance, and common motion; they constrain little about the consequences of an action outside the data.
An interventional test, which fixes the initial state, varies the action, and checks whether predicted outcomes vary as real outcomes do, is the most direct way to distinguish a model of what tends to happen from a model of what an action causes.

Accurate geometry does not imply an actionable physical model.
Mass, friction, articulation, and compliance are often visually unobservable yet critical for manipulation (Sections~\ref{subsec:p3_s7_1_2} and~\ref{sec:p4_s9_1}), and a Gaussian-splat replica of a room can be photometrically faithful while encoding nothing about which drawers open or how much force a door requires (Section~\ref{sec:p4_s11_1}).
Real-to-sim pipelines estimate such quantities by manipulating objects and observing the response \cite{pfaff2025scalable}, thereby estimating properties needed for prediction, and world-action models that predict compact encodings of future observations rather than pixels avoid spending capacity on appearance that control does not use \cite{yu2026jam}.
Confidence sets calibrated under one data-collection policy may become unreliable under policy shifts \cite{mittal2026calibrated}, and model-based action selection may induce policy shifts that invalidate this calibration.

A tractable form of the question is where, rather than whether, a world model can be trusted.
Treated like a measuring instrument, a world model would be calibrated against the environment it claims to simulate, carry a stated range of states and actions within which it is valid, and be used with caution outside it.
Establishing that range requires tests under controlled interventions and unfamiliar actions, under environmental change, and against downstream decisions, in place of average visual or predictive accuracy.
World models expressed as code, whose rules are explicit and correctable when they conflict with observation, admit such tests more readily than learned ones (Section~\ref{subsubsec:p3_s8_6_code}), and the natural sciences already hold generative models to this standard, counting a proposed material only once it has been synthesized (Section~\ref{sec:p4_s12_2}).
When an action falls outside a model's established range, the agent must decide how to obtain further evidence, which actions are safe meanwhile, and how to use the outcome to extend the model.


\section{Situated and Interactive Spatial Intelligence}
\label{sec:p5_agency}
\label{sec:p5_competence}
\label{sec:p5_people}

\FramedInitial{S}{patial} competence is not exhausted by interpreting observations or predicting future states.
In an unfamiliar environment an agent must decide where to look, what to investigate, and when to act under uncertainty, and models that perform well on provided evidence can perform poorly when responsible for gathering it.
Theory of Space makes the comparison directly: foundation models asked to explore until they knew where objects were did worse when they collected the observations themselves than when the same observations were supplied, and explored less systematically than simple scripted strategies \cite{zhang2026theory}.
On VABench, models that judge spatial relations well still fail many manipulation tasks and almost all long-horizon ones, and the gap narrows when they may move the camera and look again \cite{zhang2026vabench}.
These results locate part of the weakness in deciding what to do, rather than only in recognizing what is seen.

Section~\ref{sec:p5_observability} asked which properties can be made identifiable by acting; the question here is which action to choose.
Information gathering should be guided not only by uncertainty but by the consequences of that uncertainty for the decision at hand.
An observation that sharpens a reconstruction may be worth less than one that settles whether a passage is traversable or a grasp is safe, and the value of a view depends on the goals, costs, and risks of the actions that follow it.
Exploration policies learned in simulation formalize this for navigation objectives (Section~\ref{subsec:p3_s7_2_5}); the counterpart for the richer states of Sections~\ref{sec:p5_sufficiency} and~\ref{sec:p5_persistence}, where the quantity to resolve may be an object's identity, a person's belief, or a door's state, is largely untrained.
Reporting accuracy as a function of an interaction budget, the number of views, moves, or queries a model may spend, would place passive and active evaluation on one scale and reward agents that reach correct decisions with few observations.

Passive data, simulation, and real-world interaction provide complementary forms of experience.
Internet video, the largest source, supports perceptual and physical priors and already enables robot planning after modest fine-tuning (Section~\ref{sec:p3_s8_6}); simulation and generated worlds permit controlled exploration without risk (Section~\ref{subsec:p3_s7_2_5}); real interaction alone exposes the consequences of the agent's own decisions, which passive video does not contain and on which both the update policy of Section~\ref{sec:p5_persistence} and the predictor of Section~\ref{sec:p5_validity} depend.
Each source carries its own distortions.
The simulation gap is widest where contact, materials, lighting, and sensor noise matter (Sections~\ref{sec:p4_s9_1} and~\ref{sec:p4_s11_3}), and experience generated inside a world model is useful only within the model's range of validity.
Learning inside a world model, with rollbacks to earlier decision points and correction of the model from real transitions \cite{luo2026rollback}, lowers the cost of real interaction without removing the need for it.
How to combine the three sources without transferring incorrect assumptions from simulated or observational distributions into real-world action remains unresolved, and existing robot datasets, concentrated on tabletop success, offer little of the failures and recoveries such a combination requires (Section~\ref{subsec:p3_s7_3_2}).
The consequences of an action, moreover, are not confined to the agent's own observations: in a shared space, acting also changes what other agents observe, believe, and do.

Different agents observe and understand the same environment differently, and a shared physical world does not imply shared knowledge of it.
A system must relate its own reference frame to those of others, since a location on the user's left may lie on the robot's right and people switch between egocentric and allocentric descriptions without noticing (Section~\ref{subsec:p3_s7_3_1}).
It must distinguish the physical state from what a person believes about it, since a user who left the kitchen before the kettle was moved still believes it is on the stove.
It must anticipate how its own actions change what others perceive and do.
Mental World Modeling maintains joint physical and mental states for this reason and finds both necessary to predict human decisions (Section~\ref{subsubsec:p3_s8_6_mwm}) \cite{fei2026mental}.
People, on this view, belong in the world state on the same footing as walls and doors, with positions, viewpoints, and beliefs treated as quantities to be estimated.

In human environments, spatial actions are inherently interactive.
Robots that model people only as moving obstacles freeze in dense crowds once every path appears blocked, unable to anticipate that people will make room (Section~\ref{sec:p4_s13_1}), and an assistant's motion, speech, and help change what the people around it do next, so a plan that ignores this is planning against a stale state.
Persistent spatial memory carries a second consequence: a world state that includes people records where they live, what they own, and how they spend their days, and the design choices of what to store, at what detail, and for how long matter more than the access controls applied afterward (Sections~\ref{sec:p4_s10_1} and~\ref{sec:p4_s13_1}).
Such systems also act on their users' own spatial abilities, which weaken under habitual reliance on navigation aids and strengthen with practice \cite{dahmani2020habitual}.
Reliable spatial agency must therefore account for other agents' responses and for the consequences of maintaining and using spatial knowledge about them.

Agents that acquire the information they need, act despite incomplete knowledge, and revise their world models from the consequences of interaction would address each of the preceding questions in operational form: action changes what is identifiable, exposes which information the representation lacked, supplies the observations that keep the state current, and tests the predictions the model made.
Spatial intelligence, on this view, is a closed loop in which perception, representation, prediction, and action continually constrain and improve one another.
A usable model of physical space is one that remains open to new evidence, explicit about its limits, and corrected by what its own actions reveal.

\chapter{Research Gaps and Future Directions}
\label{chapter:p5_challenges}


% \lettrine[lines=3]{\initfamily\textcolor{darkgreen}{F}}{uture}
% Vision-based ranger involves leveraging geometric principles to calculate distance.
% A vision-based ranger could use one or multiple cameras.
% The monocular camera ranger requires prior knowledge (\eg the size of a reference object), and can only estimating in a rough distance.

\FramedInitial{T}{he} problems in this chapter are widely acknowledged and concern missing methods, inadequate measurements, and systems that work in the laboratory but not yet in open environments.
Section~\ref{sec:p5_gaps} describes research gaps along the loop of Figure~\ref{fig:p1_s2_1_intro}, Section~\ref{sec:p5_benchmarks} gaps in how progress is measured, Section~\ref{sec:p5_deployment} bottlenecks in deployment, and Section~\ref{sec:p5_directions} directions for future work.

Chapter~\ref{chapter:p5_discussions} asked why these problems are hard in principle; this chapter asks where current systems fail, how the failures can be measured, and which methods could be tested against them.
Four themes are its own: competence at tabletop or room scale does not transfer to buildings and cities; correct spatial judgments do not reliably become completed long-horizon actions; much evaluation does not measure the abilities it is named for; and deployed systems face latency, distribution shift, security, data scarcity, and multi-agent synchronization that research systems rarely confront.
Where a failure has a principled explanation in the preceding chapter, we point to it rather than repeat it.

\section{Research Gaps along the Loop}
\label{sec:p5_gaps}

\FramedInitial{M}{ost} of the following gaps are invisible in single-step evaluations and appear only when a system must integrate many views, revise its memory, act over many steps, or operate beyond room scale.

\paragraph{Models recognize scenes without relating views to one another.}
Multimodal models can name nearly every object in a room yet cannot reliably say where the objects are; Chapter~\ref{chapter:3:sec:6} found them strong on contextual and functional relations and weak on distances, sizes, and angles (Section~\ref{sec:spatial_relation_extraction}).
Given several views of a room, they rarely notice that one was captured after an object moved \cite{cao2026samespace}, and over routes spanning several streets they answer by recognizing landmarks and matching text, with little evidence of path integration \cite{mahendran2026kilometervision}.
These failures admit more than one explanation: the model may hold no account of the scene against which a later view could be checked, or it may hold spatial information that it cannot associate across viewpoints, update, or use in a metric computation.
Current benchmarks do not separate these cases (Section~\ref{sec:p5_benchmarks}).

SLAM registers each new view against a map and treats disagreement as a signal to correct one or the other (Section~\ref{sec:mapping_slam}).
Whether registering views in a common frame before language is involved (Section~\ref{cross_modality_vision_vision}), for instance with feed-forward models that regress point maps from image collections (Section~\ref{sec:visual_odometry}), closes the gap for multimodal models is unsettled, and metric estimates must in any case survive unknown or wrong camera calibration, the main open problem even for specialist pipelines (Section~\ref{sec:height_distance}).

\paragraph{Memory breaks when it has to be revised.}
Section~\ref{sec:p5_persistence} argued that persistence requires revising beliefs, not only retaining observations.
The gap is that few learned systems implement a write policy deciding what enters the world state, with what confidence, and when an entry is revised or removed.
Work on spatial memory has concentrated on recall, through longer contexts, better retrieval, and richer memory structures (Section~\ref{subsec:p3_s7_3_6}), whereas the reported failures concern revision: outdated beliefs kept after the world changes \cite{zhang2026theory}, append-only memories that accumulate contradictions, and agents that drift as their context fills (Section~\ref{subsec:p3_s7_3_7}).
Applications make the requirement concrete: a manipulator must keep an object's pose current while its own gripper occludes it (Section~\ref{sec:p4_s9_1}), and a home assistant must track object locations over weeks as the household moves things (Section~\ref{sec:p4_s13_1}).
SLAM makes a version of these decisions when it selects keyframes, merges landmarks, and discards map regions that no longer match (Section~\ref{sec:mapping_slam}); we are not aware of a learned counterpart for the richer world states that applications now require.

\paragraph{Pixel prediction is mismatched with the needs of control.}
Video world models predict pixels, while a controller needs the quantities that decide whether an action succeeds, such as contact, mass, and friction, which are often invisible in images (Section~\ref{subsec:p3_s7_1_2}).
A pixel loss spends most of a model's capacity on lighting and texture, pixel predictions drift over long horizons (Section~\ref{subsec:p3_s8_3_6}), and a model can score well on tasks with incorrect dynamics \cite{gonzalezlaiz2026limits}.
The objective rewards appearance; the use requires outcomes.
World-action models that predict compact encodings of future observations match generative ones at much lower cost \cite{yu2026jam}, following the predictive-embedding principle of the JEPA family (Section~\ref{subsec:p3_s8_6_rep}), and a further step is to predict physical quantities directly, estimated through pushing, lifting, or probing.
Whether such predictions hold under intervention is the question of Section~\ref{sec:p5_validity}; the gap here is that these quantities are rarely prediction targets at all.

\paragraph{Correct perception does not carry through to completed actions.}
Chapter~\ref{sec:reasoning} described the gap between spatial judgment and action as instructions interpreted correctly and executed incorrectly (Section~\ref{subsec:p3_s7_3_2}).
On VABench, models that judge spatial relations well still fail many manipulation tasks and almost all long-horizon ones, and degrade markedly on unfamiliar shapes \cite{zhang2026vabench}; generalist robot policies improve rapidly, yet their spatial representations drift over long action sequences (Section~\ref{subsec:p3_s7_3_2}).
Small errors compound over many steps, and in contact-rich tasks a single misjudged grasp can end an episode, so the gap is one of execution rather than judgment and grows with the horizon.
An agent must keep observing and correcting until the task is complete, as the gains from letting a model look again suggest (Section~\ref{sec:p5_agency}); this is the third open tension of Chapter~\ref{chapter:structure}, the integration of spatial reasoning into long-horizon embodied agents.
Closing it requires training and evaluating the act--observe--correct loop as a whole, with benchmarks that reward recovery as well as first-attempt success.

\paragraph{Room-scale results say little about larger spaces.}
Most results in this handbook concern single rooms or tabletops.
At street and city scale, models revert to landmark recognition even over a single kilometer \cite{mahendran2026kilometervision}, geospatial models trained on one region or sensor degrade on another (Section~\ref{sec:p4_s11_2}), aerial and ground robots must align views from very different heights (Section~\ref{sec:p4_s9_4}), and skills learned on one robot body transfer poorly to another (Section~\ref{sec:p4_s9_1}).
Scale changes more than the size of the numbers: the dominant representation shifts from local metric geometry to routes and topology and then to georeferenced coordinates, the sensing from egocentric views to aerial imagery, most of the state from seen to remembered or inferred, and the rate of change from seconds to years.

Geographic information systems handle scale by nesting reference frames, placing an object in a room, the room in a building, and the building in a georeferenced map, each level with its own coordinates and precision.
Learned spatial models rarely represent such nesting.
We expect nested frames to be part of the solution, together with benchmarks that pose the same question at several scales (Section~\ref{sec:p5_benchmarks}), and we regard cross-scale transfer as a challenge in its own right.

Each of these gaps appears where one part of the loop depends on another's output, which is why evaluating each part on its own misses them.




\section{Benchmarking Gaps}
\label{sec:p5_benchmarks}

\FramedInitial{T}{he} preface asked how the field can measure genuine progress without mistaking benchmark saturation for it.
Chapter~\ref{chapter:p5_discussions} described the properties a usable spatial model should have; this section asks how to measure them.
Most benchmarks compared in Section~\ref{sec:p3_s7_5} pose questions about a single image or video, and few require the model to act.
This format admits shortcuts, which models have repeatedly exploited: models that nearly solved synthetic visual reasoning datasets failed on novel combinations of familiar elements and on real images (Section~\ref{subsec:p3_s7_2_1}), and much of a referring-expression score survives when the words of the expression are shuffled (Section~\ref{sec:vl_grounding}).

Measured against the loop of Chapter~\ref{chapter:definition}, current evaluation covers some components far better than others.
Perception is evaluated frame by frame, through detection, segmentation, depth, and tracking accuracy, and rarely by whether its uncertainty is calibrated or its output stable over time.
Reasoning is evaluated by answer accuracy on static questions, rarely by whether the answer rests on a correct belief about the scene.
Prediction is evaluated by visual quality, rarely by correctness under actions absent from training.
Action is evaluated by task success in simulation, rarely by recovery after failure.
The update $\Psi_{\mathcal{S}}$ is rarely evaluated on its own, and neither is the effect a system has on the people who use it.

\paragraph{Score the belief behind the answer.}
A correct answer may come from language priors, regularities in the answer format, or a correct model of the scene, and accuracy alone cannot tell them apart, nor separate the explanations left open for the first gap of Section~\ref{sec:p5_gaps}.
Benchmarks can instead read out a model's belief about the scene and score it directly, as Theory of Space does at each step of an exploration \cite{zhang2026theory}, or probe hidden states, as was done for the Othello model of Section~\ref{sec:rep_language}.
A cheaper variant asks several questions about the same scene and checks that the answers are consistent with a single layout.
Controls that remove the image or shuffle the input are inexpensive and should be standard.

\paragraph{Gathering evidence is a separate ability.}
Static benchmarks supply the evidence; deployed systems must collect it.
Deciding where to look, when to move, and when to ask is a skill in its own right, and the active--passive gap of Section~\ref{sec:p5_agency} shows that current models lack it.
We propose reporting accuracy as a function of an interaction budget, the number of views, moves, or queries a model may use, which places passive and active evaluation on one scale and rewards models that reach correct answers with little effort.
In deployment the budget has a direct cost in time and energy.

\paragraph{Spatial failures unfold over time.}
Most benchmark questions refer to a single moment, while many failures in Section~\ref{sec:p5_gaps} emerge only over time.
Streaming protocols show the model only the video seen so far, so it cannot use future frames to answer about the past \cite{li2026ovosbench}; revisit protocols move the camera away and back to check whether a place is rendered consistently or changed correctly \cite{wang2026worldprobe}; change protocols alter the scene behind the agent and test whether its belief follows.
None is yet standard, and long-horizon tasks remain rare even among embodied benchmarks (Section~\ref{sec:p3_s7_5}).

\paragraph{The same question should be asked at several scales.}
Benchmarks cluster at room scale and are often built on a few indoor scan collections, while surgery, orbital operations, and underwater robotics have benchmarks of their own that rarely share protocols with them (Sections~\ref{sec:p4_s12_1} and~\ref{sec:p4_s11_3} and Chapter~\ref{sec:p4_s14}).
Posing the same spatial question at the scale of a table, a building, and a city would show where competence stops transferring and whether a model represents nested frames; a model that answers well about a desk and poorly about a building has learned the regularities of tabletop scenes, not a general spatial ability.

\paragraph{Generated worlds should be judged by what an agent can do in them.}
Generated and reconstructed worlds are still evaluated mainly by appearance, which Chapter~\ref{sec:generation} showed to be misleading.
Video metrics do not register physically impossible motion (Section~\ref{subsec:p3_s8_3_6}), rendered-view similarity cannot distinguish a solid object from a cut-out that matches the training views (Section~\ref{subsec:p3_s8_2_1}), no standard protocol exists for the physical correctness or long-term stability of world models (Section~\ref{subsec:p3_s8_6_discussion}), and vision-language judges share the blind spots of the models they evaluate (Section~\ref{sec:p3_s8_8}).
A more informative test asks whether an agent that plans or trains in the generated world succeeds in the real one, complemented by the interventional tests of Section~\ref{sec:p5_validity}.

\paragraph{Benchmarks should borrow construct validity from psychometrics.}
Spatial benchmarks adopted tasks from psychometrics, from mental rotation to perspective taking (Chapter~\ref{chapter:definition}), but not its standards.
A psychometric test reports the distribution of human performance, establishes reliability across repeated testing, and checks construct validity through agreement with other tests of the same ability and divergence from tests of different ones; most AI benchmarks report at most a single human score.
A benchmark whose questions can be answered from text alone does not measure spatial ability, and automatically generated benchmarks, now common, need such checks all the more.




\section{Deployment Bottlenecks}
\label{sec:p5_deployment}

\FramedInitial{M}{ost} spatial systems that operate reliably in deployment do so by restricting the world they must handle, a pattern that recurs across Part~\ref{part:applications}.
Driving systems are approved for a restricted operational design domain and add sensors and fallbacks to stay within it (Section~\ref{sec:p4_s9_2}); warehouse robots work in sites with fixed item sets and layouts (Section~\ref{sec:p4_s9_5}); surgical robots are mostly teleoperated and planetary rovers act within ground-supervised workflows (Sections~\ref{sec:p4_s12_1} and~\ref{sec:p4_s11_3}); headsets track most reliably in previously mapped rooms, and the first home robots rely on remote operators for tasks they cannot yet perform (Sections~\ref{sec:p4_s10_1} and~\ref{sec:p4_s13_1}).
When the domain is widened, the world state is typically the first component to fail.

\paragraph{Guarantees end at the edge of the model.}
A constraint-based planner certifies collision avoidance only for the obstacles it knows about, and a driving system is validated only on anticipated scenarios (Sections~\ref{sec:p4_s9_1} and~\ref{sec:p4_s9_2}).
Learning moves these boundaries without notice: a model accurate on average can be confidently wrong in rare situations, its uncertainty estimates can lose validity when its use changes (Section~\ref{sec:p5_validity}), and an update that improves one behavior can silently degrade another (Section~\ref{sec:p4_s9_2}).
No practical method yet states and verifies a property of a learned spatial system over a realistic domain.
Legal responsibility is bounded the same way, since under conditional automation the manufacturer is responsible for driving only within a narrow domain (Section~\ref{sec:p4_s9_2}), and each extension reopens the question.

\paragraph{The large model cannot sit in the inner loop.}
A headset must respond to head motion within tens of milliseconds, a robot arm must update its state estimate at control rate, a drone must detect small objects in real time on onboard hardware, and a spacecraft must do similar work on radiation-hardened processors far weaker than ground computers (Sections~\ref{sec:p4_s10_1}, \ref{sec:p4_s9_1}, \ref{sec:p4_s9_4}, and~\ref{sec:p4_s11_3}).
Large foundation models are much slower, even after work that accelerates them by predicting geometry internally and compressing spatial tokens \cite{li2026geothinkerf}.
A practical architecture lets a small, fast component keep the world state current at sensor rate and lets the large model consult it when time permits, mirroring the fast--slow division that Section~\ref{subsec:p3_s7_3_5} described for tool-augmented reasoning.
Whether that state is explicit, learned, or mixed is the question of Section~\ref{sec:p5_persistence}; the deployment constraint is only that it be maintained at sensor rate and remain consultable.

\paragraph{A map starts going stale as soon as it is made.}
Furniture is moved, shop displays change, buildings are renovated, crops grow, and tissue shifts during surgery, so a deployed map must be updated continually, which Part~\ref{part:applications} found difficult in every domain.
How a state should be revised is the algorithmic question of Section~\ref{sec:p5_persistence}; the deployment problem is the cost and latency of keeping it revised and the synchronization of its copies across devices.
Digital twins call the discrepancy between model and asset the liveness gap and treat closing it at scale as open (Section~\ref{sec:p4_s11_1}); lifelong SLAM must withstand months of appearance change (Section~\ref{sec:mapping_slam}); shared augmented-reality maps are hard to keep current (Section~\ref{sec:entertainment}); surgical navigation loses its preoperative registration once the brain shifts (Section~\ref{sec:p4_s12_1}).
We suggest that deployed systems report how far their world state lags behind the world it describes, a quantity that is rarely measured.

\paragraph{Real data are scarcest where simulation is least trustworthy.}
Spacecraft perception is trained mostly on simulated images, and simulators do not reproduce the lighting of orbit, reflections from spacecraft materials, or radiation noise (Section~\ref{sec:p4_s11_3}); touch, force, and sound have the least data of all modalities and are also hard to simulate (Sections~\ref{subsec:p3_s7_3_8} and~\ref{sec:p4_s9_1}); clinical data are scarce because of privacy regulation and the cost of expert annotation (Section~\ref{sec:p4_s12_1}).
In the first two cases simulation cannot compensate, since the simulator is least accurate for exactly the effects the missing data would capture.
Simulations built from real measurements, data engines that target rare cases, and cross-institution training without sharing raw data all help, but none yet operates at the breadth that general spatial competence requires.

\paragraph{A persistent map is also an attack surface.}
Geospatial data raise issues of sovereignty and adversarial manipulation (Section~\ref{sec:p4_s11_2}), a compromised digital-twin server exposes the layout and vulnerabilities of the asset it mirrors (Section~\ref{sec:p4_s11_1}), and world-action models can be attacked through their visual inputs, with small perturbations producing large changes in behavior \cite{wu2026pasakv}.
These risks grow with the persistence of a map and the number of parties sharing it, turning the privacy concerns of Section~\ref{sec:p5_agency} into security concerns.




\section{Future Directions}
\label{sec:p5_directions}

\FramedInitial{T}{he} following directions build on the gaps above and on the questions left open in Chapter~\ref{chapter:p5_discussions}.
Each is a hypothesis rather than a recommendation: where it takes an architectural position, we name the alternative it must beat and an experiment that would decide.

\paragraph{Persistence and rendering can be separated.}
Interactive video models remember a place by regenerating it, which ties memory to the context window and to the renderer.
A factorized design would keep what must persist, such as layout, object identities, and their uncertainty, in an explicit layer, and use a learned layer to predict change and render what the agent needs to see, with $\Psi_{\mathcal{S}}$ promoting confirmed predictions into the explicit layer and revising it when observations contradict it.
The alternatives are a latent state with learned memory, as in models that factorize memory inside the generator \cite{yu2026grome}, and a growing context with retrieval.
The hypothesis is that at equal compute the factorized design stays consistent over longer horizons and more changes; it is refuted if a latent-only model matches it on revisit and change protocols while passing the faithfulness tests of Section~\ref{sec:p5_validity}.
The decisive experiment is an agent that leaves a building, returns hours later, detects what changed, and acts accordingly; for the factorized variant, the open question is the interface that decides when a prediction enters the explicit layer and when an entry is overridden.

\paragraph{The medium of spatial reasoning can be learned per query.}
Section~\ref{sec:p5_sufficiency} found that spatial answers improve when a model can work with maps, imagined views, or geometric tools.
The choice of medium can be learned per query from whether it led to a correct answer and at what cost, so that simple questions are answered in language and a map or simulation is invoked only when needed, with the intermediate steps later distilled into the model (Section~\ref{subsec:p3_s7_3_8}).
The comparison is with the same model reasoning in text alone and with a policy that always invokes the tool, on map building and perspective taking, with image-removed controls; the hypothesis fails if the learned selector does not beat the always-tool policy on cost at equal accuracy.

\paragraph{The update policy can be trained from experience.}
Passive data leave the update policy of Section~\ref{sec:p5_persistence} largely untrained (Section~\ref{sec:p5_agency}).
It can be trained directly, by rewarding observations that add information to the map and penalizing beliefs kept after the evidence has changed, with learning inside a world model lowering the cost by trying alternatives in imagination first (Section~\ref{subsec:p3_s7_3_7}); benchmarks that read out a model's belief (Section~\ref{sec:p5_benchmarks}) supply the training signal as well as the evaluation.
The target is an agent that builds as accurate a world state from its own exploration as from a curated set of views, and the scripted strategies that currently outperform foundation models \cite{zhang2026theory} are the baseline to beat.

\paragraph{Verification is feasible where part of the state is explicit.}
Neural networks are difficult to verify, whereas objects, poses, and programs can be checked against constraints.
The Algebraic World Model turns exploration into hypotheses checked by machine-verified proofs \cite{pan2026algebraic}, and uncertainty can be verified once its calibration is matched to the policy actually executed \cite{mittal2026calibrated}.
The alternative is to verify latent states indirectly, through probes and faithfulness tests \cite{gonzalezlaiz2026limits}, which is cheaper and yields statistical rather than logical guarantees.
In either case the obstacle is perception, since any guarantee about the state inherits the errors of its estimate from sensor data.
A machine-checked property, such as a bound on collision probability, holding for a learned agent over a deployment-scale domain would change the situation of Section~\ref{sec:p5_deployment}; no current system provides one.

\paragraph{A shared world state needs the machinery of version control.}
When several agents and people maintain one world state, they need a common reference frame, stable object identities, a record of changes, a procedure for resolving conflicts, and permissions on who may read and write what.
Version control solved the analogous problems for shared code, and multi-agent systems that build cognitive maps already record changes as atomic commits and cross-check one another \cite{zhang2026cocosi}.
Embodied teams add difficulty, since ground and aerial robots must first align very different viewpoints (Section~\ref{sec:p4_s9_4}) and city-scale and planetary digital twins supply the georeferenced frames into which local maps must fit (Sections~\ref{sec:p4_s11_1} and~\ref{sec:p4_s11_2}); permissions are also where the privacy requirements of Section~\ref{sec:p5_agency} can be enforced.
A suitable test is a team of heterogeneous robots and a person who maintain one consistent world state of a building over several days, each able to query how the others currently perceive it.

\paragraph{A theory of spatial representation has a concrete agenda.}
Section~\ref{sec:p5_sufficiency} proposed sufficiency relative to tasks and actions as the criterion for a representation.
Turning it into a theory requires identifiability results stating which world-state variables are recoverable from a given set of modalities and actions, extending the sensor-specific analyses of Chapter~\ref{chapter:3:sec:6}; sufficiency results stating which variables must be preserved for a task family under every action the agent may take, of which equivariance to rigid motion (Section~\ref{subsec:p3_s7_2_3}) is the one case already worked out; faithfulness tests deciding, for any architecture, whether a latent state's dynamics follow the environment's \cite{gonzalezlaiz2026limits}; and a specification of which parts of a representation must be explicit for inspection and verification.
Such results would let the world-centered property required by Chapter~\ref{chapter:definition} be tested for any system, regardless of architecture.

\chapter{Conclusion}
\label{chapter:p5_conclusion}


% \lettrine[lines=3]{\initfamily\textcolor{darkgreen}{C}}{onclusion}
% Vision-based ranger involves leveraging geometric principles to calculate distance.
% A vision-based ranger could use one or multiple cameras.
% The monocular camera ranger requires prior knowledge (\eg the size of a reference object), and can only estimating in a rough distance.

\FramedInitial{P}{oincar\'e} wrote in 1902 that to localize an object is to represent to ourselves the movements needed to reach it, and this handbook opened with that sentence.
The work surveyed here supports his view, as the systems that act most reliably are those whose representations are most closely tied to prediction and action, and their main weakness is keeping those representations consistent over time.

This handbook set out to provide a common map of a fragmented field, reading work from SLAM and computer vision to geographic information science and the cognitive science of space through the representations each of them constructs.
Part~I traced five eras of spatial representation and proposed a definition in which a system and its environment are coupled through a world-centered representation that perception writes, action reads, reasoning and prediction operate on, and the update $\Psi_{\mathcal{S}}$ keeps current.
Parts~II to~IV followed these functions to the state of the art.
In the terms of Chapter~\ref{chapter:definition}, different communities built different parts of this loop: vision and graphics the representation, robotics the operators that act on it, SLAM an update with geometric guarantees; few systems close the entire loop.

Part~V located the open problems mainly at the junctions between these parts.
The grocery-handling robot introduced at the beginning of this part fails when its perception, memory, prediction, and action cease to agree about the kitchen, even though each capability works in isolation.
Chapter~\ref{chapter:p5_discussions} asked what such a shared state must provide to be usable: what can be identified from limited evidence, what must be represented, how beliefs are kept current, when predictions can be trusted, and how knowledge is acquired through interaction in space shared with others.
Chapter~\ref{chapter:p5_challenges} showed the same state failing to persist in research systems, rarely measured by benchmarks, and kept manageable in deployment only by restricting the world, and stated directions for closing these gaps as hypotheses that can be tested.

The preface posed four questions without claiming final answers, and Part~V suggests provisional ones.
A system is spatially intelligent, beyond its performance on particular tasks, when it maintains a world state that stays consistent as it moves and acts and that can serve tasks it was not designed for.
Geometric, semantic, and commonsense structure integrate best when each part of the state is stored according to how long it must persist and whether it must be inspected, and modalities as different as pixels, point clouds, language, and audio contribute as complementary constraints on one underdetermined state, each resolving ambiguities the others leave open; registering views in a common frame before language is involved appears to be the most promising route to cross-view alignment, though whether it closes the gap remains to be shown.
Spatial knowledge is most likely to transfer between environments when it is anchored in nested frames, from objects to rooms to maps.
Planning and action emerge from spatial understanding through prediction that is tested under intervention and through a loop in which the agent keeps observing and correcting.
Genuine progress shows in how well a system gathers evidence, keeps its beliefs current, and uses them, more than in its answers to static questions.

Chapter~\ref{chapter:definition} called its definition a definition in motion, and Part~V suggests three revisions to it.
The definition refers to a single system $\mathcal{S}$, whereas deployed systems increasingly share spaces and world states with other agents and with people, so it requires a plural form.
It requires the representation to support perception, reasoning, prediction, and action, and deployment adds the requirement that planners, regulators, and the people it describes can inspect it.
It admits social settings only as a property of the environment, whereas Chapter~\ref{chapter:p5_discussions} argued that people belong in the state itself, with their positions, viewpoints, and beliefs predicted and respected alongside walls and doors.

As the preface anticipated, some of the systems discussed in this handbook will soon be superseded, and some of the positions taken here will need revision as new evidence arrives.
We hope that a clear account of the open problems makes that revision easier, and that the next account of spatial intelligence finds more of the loop closed.

